\section{検定 -- 仮説と\texorpdfstring{$p$}{p}値をどう読むか}

  \subsection{この差は偶然か}

    前章の終わりに, 答えを保留した問いが一つ残っている.
    配信サービスの邦画作品の標本で, 原作のある作品（既存作品の映画化）と原作のない作品（オリジナル脚本）の平均視聴回数を比べたところ, 原作あり群のほうが多かった.
    ところが, 差に幅をつけてみると, 95\%信頼区間の下端はゼロのすぐ近くまで来ていた.
    標本を取り直していたら, ゼロをまたぐ区間が出ていてもおかしくない.
    この差を根拠に「原作の有無は視聴回数に効いている」と言ってよいのか.
    それとも, たまたま視聴回数の多い作品が原作あり側に偏っただけなのか.

    信頼区間からわかるのは, 差の見積もりがどのくらいの幅を持つかまでだ.
    目の前の差は, 偶然のばらつきとして片づけられる範囲に収まっているのか, それを超えているのか.
    この一歩踏み込んだ問いに白黒をつける手続きが, 前章の終わりに名前だけ挙がった\textbf{統計的検定}だ.

    同じ形の判断は, すでに身の回りのあちこちで使われている.
    選挙の開票速報で, 開票率がまだ数\%なのに「当選確実」が出る場面を見たことがあるだろう.
    あれは, 出口調査の標本に表れた候補者間の得票差が, 偶然では説明しにくいほど大きいかどうかを判断した結果だ.
    IT企業がWebサイトのボタンの色を変えてクリック率を比較する, いわゆるA/Bテストも同じ構造をしている.
    どの場面でも, 突き詰めれば二つのグループの数字を見比べる話である.
    見比べるだけなら, 平均を並べれば済みそうだ.
    まずはそこから試してみる.


  \subsection{目で見た差は当てになるか}

    二つのグループの平均を計算して, 差を見る.
    差が大きければ「違いがある」, 小さければ「違いはない」.
    日常の多くの場面で, 実際にこうやって判断している.

    たとえば, 二つのクラスの英語テストの平均点を比べるとしよう.
    A組の平均は72点, B組の平均は68点で, 4点の差がある.
    「けっこう違う. 教え方に差があったのではないか」と感じる人もいるだろう.

    ここで, 前章の計算を持ち込んでみる.
    各クラスの人数が15人で, 標準偏差が12点だったとする.
    前章のケース3で使った「差の標準誤差」を計算して, 二つのグループの平均の差が偶然でどのくらいばらつきうるかを見積もる.
    %
    \[
      \mathrm{SE}_{\text{差}} = \sqrt{\frac{12^2}{15} + \frac{12^2}{15}} = \sqrt{9.6 + 9.6} \approx 4.4 \text{ 点}
    \]
    %
    根号の中の各項は, 各グループの分散を人数で割った$s^2 / n$で, それぞれ$12^2 / 15 = 9.6$だ.
    前章で見たとおり, 二つのグループが独立なら, 差のばらつき（分散）はそれぞれのばらつきの足し算になる.
    結果はおよそ4.4点.
    観測された差の4点より, 偶然で見込まれるばらつきのほうが大きい.
    二つのクラスの母平均がまったく同じだったとしても, 標本平均の差は4点程度まで容易にばらつく, ということだ.

    平均点だけ見れば「けっこう違う」ように映る72点と68点も, 一人ひとりの点数が12点もばらつく中では, 偶然の差と区別がつかない.
    \textbf{差の大きさだけでは, 意味のある差かどうかは判断できない.}
    ばらつきと人数を織り込んだ, 判断の組み立てが要る.


  \subsection{「差がない世界」を出発点にする}

    二つのグループに差があるかどうかを調べたい.
    素朴に考えれば, 「差がある」ことを直接示しにいけばよさそうだ.
    ところが, ひとくちに「差がある」と言っても, その中身は一つに決まらない.
    2点の差も「差がある」だし, 10点でも0.001点でも「差がある」.
    「差がある世界」は無数にあって, 議論の出発点をどこに置けばよいのかが定まらないのだ.

    一方, 「差がゼロ」の世界はただ一つに決まる.
    そこで検定では, 発想を逆転させる.
    \textbf{「差がない世界」を仮に設定する}ことから始めるのだ.
    この世界では, 二つのグループの母平均は等しく, 観察される差はすべて偶然のばらつきから生じる.
    その世界を基準にして, 手元のデータに出た差がどのくらい珍しいかを確率で測る.
    もし「差がない世界ではめったに起こらない」ほど珍しければ, 差がないという仮定のほうに無理があると判断する.

    この「差がない」側の仮定を\textbf{帰無仮説}（$H_0$）と呼ぶ.
    対して, 「差がある」という主張を\textbf{対立仮説}（$H_1$）と呼ぶ.
    帰無仮説という名前は物々しいが, 中身は「差がない側の仮定」以上のものではない.
    ただ, 論文や記事には頻繁に出てくる用語なので, 名前を知っておくと読むときに困らない.

    \subsubsection*{差を標準化する}

    手元のデータに出た差の珍しさを, 差がない世界を基準にして測りたい.
    ところが, 前の節で見たとおり, 4点という差が大きいのか小さいのかは, データのばらつきと人数によってまったく変わる.
    差をそのまま眺めても, 珍しさの見当がつかない.

    そこで, 差をそのばらつき（標準誤差）で割って標準化する.
    前章の信頼区間でも, $\bar{x}$と$\mu$のずれを$\mathrm{SE}$の何倍分と数えて幅を作っていた.
    検定でも同じ発想を使う.
    この標準化した値を\textbf{検定統計量}と呼ぶ.
    たとえば二つの平均の差を比べるなら,
    %
    \[
      t = \frac{\text{二つの平均の差}}{\text{差の標準誤差}}
    \]
    %
    分子は差の大きさで, 分母はその差が偶然でどのくらいばらつくかである.
    差が大きいほど, ばらつきが小さいほど, 人数が多いほど, この値は大きくなる.
    ばらつきで測り直した値は, テストの点数やクリック率といった元の単位を引きずらない.
    だから, 種類の違うデータでも, 差の珍しさを同じ基準で比べられる.

    問いの種類に応じて, 使う検定統計量は変わる.
    平均の差なら$t$統計量, カテゴリの偏りならカイ二乗統計量$\chi^2$, ばらつきの比なら$F$統計量だ.
    統計量が違っても, 「差を標準化して珍しさを測る」という構造は同じである.
    どれを使うかは問いの形で決まるので, すべてを暗記する必要はない.

    \subsubsection*{珍しさを数値にする：\texorpdfstring{$p$}{p}値}

    検定統計量が, たとえば$t = 2.3$と出たとしよう.
    この数字だけを眺めていても, 「珍しい」のか「よくある」のかは判断がつかない.

    頼りになるのは, 帰無仮説のもとで検定統計量がどんな値をとりやすいか, という分布である.
    差がない世界でも, 偶然のばらつきのせいで$t$はゼロちょうどにはならず, ゼロの周りに散らばる.
    その散らばり方は理論的にわかっていて, 平均の差の$t$統計量なら\textbf{$t$分布}, カテゴリの偏りのカイ二乗統計量ならカイ二乗分布と呼ばれる分布に従う.
    この分布の中で, いまの値と同じかそれ以上に極端な値が出る確率を求める.
    それが\textbf{$p$値}である.
    $p$値が小さいほど, 差がない世界では起こりにくい結果が出ている, と読める.

    ここで, 「$p$値が小さいなら, 帰無仮説が正しい確率が低いのだろう」と読みたくなるかもしれない.
    しかしそうではない.
    $p$値は, 帰無仮説が正しい\textbf{と仮定したうえで}計算した確率であって, 帰無仮説そのものが正しいかどうかの確率ではない.

    がん検診に同じ構造がある.
    検診で「陽性」と出ても, 実際にがんである確率は数\%にとどまることがある.
    「陽性」が意味しているのは「がんがないとしたら, この検査結果は珍しい」というところまでで, 「がんである確率が高い」こととは別の話だ.
    $p$値と帰無仮説の関係も, 同じ形をしている.

    \subsubsection*{判断の基準を先に決める}

    $p$値の「どのくらい小さければ珍しいと言えるか」を, その都度決めていたらどうなるか.
    都合のよい結果が出たときだけ, 「$p = 0.06$だけど, まあ有意と言ってもいいだろう」と基準を緩めたくなる.

    この恣意性を防ぐために, 判断の基準となる値を\textbf{有意水準}（$\alpha$）と呼び, 検定の前に決めておく.
    広く使われているのは$\alpha = 0.05$で, 差がない世界でも20回に1回は有意という結果が出てしまう程度の線にあたる.
    0.05という数字に絶対的な根拠があるわけではない.
    歴史的な慣習として定着したものだ.

    ただし, どの分野でも0.05というわけではない.
    素粒子物理学では「5$\sigma$」と呼ばれる基準が使われており, $p \approx 0.0000003$に相当する.
    ヒッグス粒子の発見で採用された基準だ.
    人命に関わる医薬品の認可でも, 通常より厳しい基準が求められることがある.
    どこに線を引くかは, 判断を間違えたときのコストに応じて人間が決めることなのだ.

    \begin{mybox}{新薬の認可と二つの誤り}
      新薬の治験では, 二つの誤りが問題になる.
      一つは, 効かない薬を「効く」と判断してしまう誤り（\textbf{第I種の誤り}）.
      もう一つは, 効く薬を「効かない」と判断してしまう誤り（\textbf{第II種の誤り}）.
      第I種の誤りが起こる確率は, 有意水準$\alpha$で抑えられる.
      第II種の誤りは, サンプルサイズを大きくすることで減らせる.
      アメリカのFDA（食品医薬品局）の新薬認可では通常$p < 0.05$が要求されるが,
      がん治療薬のように「効く薬を承認しないこと」のコストが大きい場面では, 基準が変わることもある.
      0.05という数字そのものは慣習でも, どの水準を採るかは誤りのコストとの相談で決まっている.
    \end{mybox}

    「$p = 0.04$なら有意で, $p = 0.06$なら有意でない. 境界の前後でそんなに違うものか」と感じるかもしれない.
    もっともな疑問で, 実際, 0.04と0.06のあいだで何かが不連続に変わるわけではない.
    $p$値は連続的な指標であり, 有意水準は便宜上の区切り線にすぎない.
    $p$値だけで白黒をつけることの限界は多くの統計学者が指摘しており,
    効果の大きさやサンプルサイズとあわせて判断することが求められている.


  \subsection{検定の手順}

    ここまでの部品を並べ直すと, 検定は「差がない世界を仮に設定し, 手元のデータがその世界でどのくらい珍しいかを調べる」手続きとしてまとまる.

    ただし, どんなデータにでも使えるわけではない.
    観測どうしが独立であること（ある観測が別の観測に影響していないこと）, 比較するグループの定義が明確であること, そして有意水準を検定の前に決めてあることが前提になる.
    このうえで, 次の5つのステップを順に踏む.

    \begin{enumerate}
      \item \textbf{何を比較するかを決める.}
        調べたい問いに応じて, 適切な統計量を選ぶ.
        二つのグループの平均を比べたいなら$t$統計量,
        カテゴリの偏りを調べたいならカイ二乗統計量$\chi^2$,
        ばらつきの比を見たいなら$F$統計量だ.
        統計量の選択を間違えると, 「答えが出てこない」のではなく「的外れな問いに答えてしまう」ことになる.

      \item \textbf{帰無仮説を立てる.}
        「差がない」「効果がない」「関係がない」と仮定する.
        帰無仮説のもとで検定統計量が従う分布は理論的にわかっているので, 珍しさを確率として計算できる.

      \item \textbf{データから検定統計量を計算する.}
        実際のデータを使って, 差をばらつきで標準化した値を求める.

      \item \textbf{帰無仮説のもとでの分布から$p$値を求める.}
        帰無仮説が正しいとしたとき, 計算した統計量と同じかそれ以上に極端な値が出る確率だ.

      \item \textbf{$p$値と有意水準を比較して判断する.}
        $p$値が有意水準$\alpha$より小さければ, 帰無仮説のもとでは今のデータは珍しすぎると判断し,
        帰無仮説を\textbf{棄却}する. 対立仮説, つまり「差がある」という主張が支持される.
        $p$値が$\alpha$以上であれば, 帰無仮説を棄却できない.
    \end{enumerate}

    「帰無仮説を棄却できない」ことと「差がない」ことは同じではない.
    裁判にたとえるなら, 「証拠不十分で無罪」と「無実の証明」は異なる.
    今回のデータでは帰無仮説を退けるだけの証拠が見つからなかった, というだけだ.
    差が本当にゼロであると積極的に主張しているわけではない.
    「差がないこと」を直接示すには, 同等性検定という別の手順が必要になる.
    本書では扱わないが, 必要になったときに調べればよい.

    \subsubsection*{片側検定と両側検定}

    検定の前に決めておくのは, 有意水準だけではない.
    問いの向きも, データを見る前に決めておく必要がある.
    「原作あり群のほうが高い」と問うのか, 向きを決めずに「単に差がある」と問うのかで, $p$値の計算方法が変わるからだ.

    特定の方向にだけ関心がある場合を\textbf{片側検定}, 方向を問わない場合を\textbf{両側検定}と呼ぶ.
    片側検定では「極端」の定義が分布の片方の裾だけになるため, $p$値は両側の場合の半分になる.
    同じデータでも, 問い方が変われば結論が変わりうる.
    検定は, データから問いを探し出す手続きではなく, あらかじめ立てた問いに答える手続きなのだ.

    迷ったときは両側検定を選ぶのが安全だ.
    片側検定を使うのは, 理論的に一方向しかありえない場合や, データを見る前に方向まで決めていた場合に限る.
    分析が終わってから都合のよい方向を選ぶのは, 判断の基準を後から緩めるのと同じことになる.

    \begin{mybox}{\textit{Student}とビール工場}
      $t$統計量を考案したのは, ウィリアム・ゴセット（1876--1937）というイギリスの統計学者だ.
      彼はダブリンのギネス醸造所で品質管理を担当していた.
      少数のビールサンプルから品質を判断しなければならないという実務的な動機から,
      小標本でも使える検定法を開発した（1908年）.
      会社の規定で本名での論文発表ができなかったため, ``Student''というペンネームを使った.
      $t$分布を``Student's $t$''と呼ぶのはこのためだ.
    \end{mybox}


  \subsection{原作の有無で平均視聴回数に差はあるか}

    冒頭の問いに, 検定の手順を当てはめてみよう.

    手元の標本にある作品を, 原作の有無で二つのグループに分ける.
    一方は原作のある作品の集まり（原作あり群）,
    他方は原作のない作品の集まり（原作なし群）だ.
    このように, 別々の対象を二つのグループに分けて比べる方法を\textbf{対応のない二群の比較}と呼ぶ.

    5つのステップに沿って進める.

    何を比較するかは決まっている.
    二つのグループの平均視聴回数だから, $t$統計量を使う.
    帰無仮説は「原作の有無によって平均視聴回数に差はない」.
    対立仮説は「原作あり群のほうが平均視聴回数が高い」とする.
    方向を指定しているので, 片側検定にあたる.
    「単に差があるかどうか」だけを知りたいなら両側検定になる.
    どちらにするかは, データを見る前に決めておくのだった.

    検定統計量は, 二つのグループの平均とばらつきから$t$統計量として計算する.
    $p$値は, 「平均視聴回数に差がない」と仮定したときに, 得られた$t$統計量と同じかそれ以上に極端な値が出る確率として, $t$分布から求める.
    具体的な計算はPart Bで手を動かしながら確認する.

    あらかじめ有意水準を$\alpha = 0.05$と決めていたとしよう.
    $p$値が0.05より小さければ帰無仮説を棄却し, 「原作あり群のほうが平均視聴回数が高い」という主張が支持される.
    0.05以上であれば, 帰無仮説は棄却できない.

    ただし, 仮に有意な差が出たとしても, 「原作の有無」そのものが視聴回数を押し上げたとまでは言い切れない.
    原作のある作品は知名度や製作費, 宣伝規模が大きい傾向があり, その違いが視聴回数の差として表に出ている可能性がある.
    検定で言えるのは「平均視聴回数の差は偶然では説明しにくい」というところまでで, 何がその差を生んだのかには別の議論が要る.
    この種の入り組んだ関係は, 9章で改めて扱う.

    章の冒頭で, 差の95\%信頼区間の下端がゼロのすぐ近くまで来ていることを気にしていた.
    あのとき見ていたものは, 実は検定と同じである.
    二つの平均の差の95\%信頼区間がゼロを含んでいなければ, 有意水準5\%の両側検定で帰無仮説は棄却される.
    逆にゼロを含んでいれば, 棄却できない.
    信頼区間は差の大きさを幅で見積もり, 検定は「偶然で説明できるか」に判定を下す.
    同じ情報を, 別の角度から読んでいるのだ.
    この対応も, Part Bの計算で確かめられる.


  \subsection{血液型と性格をカイ二乗検定で調べる}

    今度は, 量的データではなくカテゴリデータの検定を見てみよう.

    一昔前に, 血液型によって性格がわかるという診断が流行したことがある.
    「A型は几帳面, B型はマイペース」といった話だ.
    こうした主張に科学的根拠があるかどうかはひとまず脇に置いて,
    「血液型と性格に関係があるのか」をデータで調べるとどうなるかを見ていく.

    100人に「血液型」と「性格タイプ（几帳面／大雑把）」を聞いた結果が, 次のクロス集計表にまとまっているとする.

    \begin{center}
      \begin{tabular}{l|cccc|c}
             & A型 & B型 & O型 & AB型 & 合計 \\
      \hline
      几帳面 & 18  & 6   & 10  & 8    & 42 \\
      大雑把 & 7   & 19  & 15  & 17   & 58 \\
      \hline
      合計   & 25  & 25  & 25  & 25   & 100
      \end{tabular}
    \end{center}

    A型で「几帳面」と答えた人は18人.
    他の血液型より明らかに多く, 「A型は几帳面」という俗説を裏づけているように見える.

    その「多い」は, 何と比べての多さだろうか.
    全体では100人中42人が「几帳面」と答えており, 各血液型は25人ずつだ.
    もし血液型と性格がまったく無関係なら, どの血液型でも几帳面な人は$25 \times 42/100 = 10.5$人くらいになるはずである.
    この10.5人を\textbf{期待度数}と呼ぶ.
    「無関係だとしたら, 行の合計と列の合計の比率から, 表の各マスにはこのくらいの人数が入るはず」という数値だ.

    A型で几帳面な人は18人で, 無関係なら10.5人程度.
    たしかにずれている.
    問題は, このずれが偶然でも起こりうるのか, それとも偶然では説明しにくいのか, である.

    手順は先ほどの$t$検定と同じ構造をたどる.
    帰無仮説は「血液型と性格タイプは無関係である」.
    各マスについて, 行合計$\times$列合計$\div$全体で期待度数を計算する.
    次に, 実際の人数と期待度数のずれを, すべてのマスについて集計してカイ二乗統計量を求める.
    ずれが大きいほど, カイ二乗統計量は大きくなる.
    最後に, 帰無仮説のもとでのカイ二乗分布から$p$値を求め, 有意水準を下回れば「無関係とは考えにくい」と判断する.
    具体的な計算はPart Bで確認する.

    カイ二乗検定の発想には100年以上の歴史がある.
    カール・ピアソンが1900年に, サイコロの出目が本当に均等かどうかを調べるために考案したのが始まりだ.

    なお, 仮にこの検定で有意な結果が出たとしても, それだけで「血液型と性格に関係がある」とは言い切れない.
    「血液型の俗説に合わせて回答してしまう人がいる」「性格の分類が恣意的である」など, データの集め方に由来する偏りが結果に影響している可能性があるからだ.
    検定にできるのは, データ上の偏りが偶然では説明しにくいと示すところまでで, その偏りが何に由来するかは検定の外の問題である.
    調査設計やバイアスの問題は9章で改めて扱う.


  \subsection{何がわかって, 何がまだ足りないか}

    出発点は, 「二つのグループに見えた差は, 偶然のばらつきで説明できるのか」という問いだった.
    目で見た差をそのまま信じると, 判断を誤ることがある.
    そこで, 差がない世界を仮に設定し, ばらつきとサンプルサイズを織り込んだうえで, 手元のデータの珍しさを測る.
    この枠組みが統計的検定だった.

    枠組みを支える部品は四つある.
    帰無仮説と対立仮説で問いを定式化し, 検定統計量で差を標準化し, $p$値で珍しさを数値に直し, 有意水準で判断の基準を事前に固定する.
    この四つがそろって初めて, 「この差は偶然で説明できる範囲か」を, 誰がやっても同じ手順で判断できる.

    ただし, 検定が答えているのは「偶然で説明できるか」までで, 「意味のある差か」には答えていない.
    ある健康食品の臨床試験を想像してほしい.
    10万人規模の調査で, 食品を摂取したグループの血圧が平均0.3\,mmHg低かったとする.
    $\mathrm{SE}$は$\sigma / \sqrt{n}$で計算するから, $n = 100{,}000$ともなると$\mathrm{SE}$は非常に小さくなり, 0.3\,mmHgの差でも$p < 0.001$となりうる.
    しかし, 0.3\,mmHgは血圧計の測定誤差の範囲にすら収まるような差で, 臨床的にはまったく意味がない.
    \textbf{「統計的に有意」であることと「実用的に意味がある」ことはまったく別}である.
    サンプルサイズが巨大なら, どんなに小さな差でも$p$値は小さくなる.
    逆にサンプルサイズが小さければ, 本当に差があっても検出できないことがある.
    $p$値が小さいことは, 「帰無仮説が間違っている」ことの証明でも, 「効果が大きい」ことの証拠でもない.
    差の大きさそのものは, 効果量という別の指標で測る.
    Part Bと9章で改めて扱う.

    2011年には, 心理学者ダリル・ベムが「未来予知」の実験で$p < 0.05$を得て一流誌に論文が掲載され, のちの追試では再現されない, という事件も起きている.
    $p$値だけで結論を出すことの危うさを象徴する事例だ.

    本章では, 検定が何をする手続きなのか, なぜそうするのかを, 数式を最小限にして追ってきた.
    次のPart Bでは, 同じ道を計算で歩き直す.
    二群の$t$検定（Welch法）とカイ二乗検定を実際に手で計算しながら,
    検定統計量がどう求まり, 自由度がどう決まり, $p$値がどう読めるのかを確かめていく.
